\subsection{Direct-evaluator cost estimates}
\label{app:computational-cost}

We tested the shortcut-free projector evaluator on deliberately elementary
categorical data, $SU(2)_2$, and on a familiar simply connected manifold,
$E(2)=K3$.  The diagram was a draft transcription of the $n=2$ specialization
of Gompf--Stipsicz Figure~8.15 \cite[Figure~8.15]{GompfStipsicz1999}.  It has
$22$ ordinary components and therefore $3^{22}=31{,}381{,}059{,}609$
simple-label sectors.  The transcription passed recorded link-level checks
but was not certified against the printed figure, and no $E(2)$ invariant was
completed.

An early informal extrapolation was roughly $350$ serial years, but we do not
use that number: no current reproducible artifact records its assumptions.
The retained version-1 runtime-floor artifact instead makes assumptions
chosen to be unrealistically favourable: sampled fast strata are extrapolated,
fast sectors cost zero, and every remaining sector costs exactly $0.01$
seconds.  Even under those assumptions the total is
$313{,}809{,}868.7$ worker-seconds, or $87{,}169.4$ worker-hours and
$9.94$ worker-years.  It gives $30.27$ days on $120$ perfectly utilized
workers.  Ideal linear scaling would require approximately $3{,}633$ workers
for one day, $87{,}170$ for one hour, or $5.23$ million for one minute.  In
the formal limit of one worker per sector, the toy model reaches $0.01$
seconds only by using approximately $31.38$ billion workers and ignoring
scheduling, communication, memory, and the hard tail.  Thus the outer sum is
embarrassingly parallel, but the contraction inside a hard sector is not
automatically so.

For a reproducible cloud-price scale, consider Hetzner's European CCX43
instance: $16$ dedicated AMD vCPUs and $64$ GB RAM.  Hetzner's price schedule
effective June~15, 2026 lists €0.4423 per hour and a €275.99 monthly cap,
excluding VAT and IPv4 \cite{HetznerCloudPricing2026}.  Identifying one
perfectly utilized worker with one dedicated vCPU converts the optimistic
work volume into $5{,}448.1$ CCX43-hours, or €2,409.69 at the uncapped hourly
rate.  A concrete $120$-worker deployment requires eight such instances
($128$ vCPUs); running them for $30.27$ days costs €2,570.34 at the hourly
rate, or approximately €2,207.92 if all eight instances receive exactly one
monthly cap.  These figures exclude VAT, IPv4, storage, excess traffic,
orchestration, failures, imperfect scaling, and the unmeasured hard tail.
They are therefore pricing illustrations for the optimistic toy model, not
quotes or upper bounds for completing the computation.

A quantum computer does not presently turn this into a defensible runtime
estimate.  If every sector contribution could be implemented as a bounded,
coherent unit-cost oracle and only an additive approximation were required,
quantum amplitude estimation could give its usual quadratic query advantage
\cite{BrassardEtAl2002}.  Our contractions do not supply such an oracle for
free, and exact summation is not covered by that statement.  In particular,
replacing $3^{22}$ by its square root $3^{11}=177{,}147$ is not a valid
runtime prediction for the present evaluator.  A quantum estimate would
first require a reversible circuit for one sector, an error target, and a
fault-tolerant gate-cost analysis.
